\documentclass[fontset=none]{ctexart}
\usepackage{qp-m3}
\begin{document}
\renewcommand{\qpzhangming}{第一章\quad 空间向量与立体几何}
\renewcommand{\qpjianming}{测评卷·试产片3}
\keshi{测评卷（卷件场·页脚灰块＋白底答案块基准）}
\noindent{\fangsong 本页举证两场：①页脚＝YJ 组灰块页码（卷件形制，页码零头挂）；②答案块＝A2 白底制式（卷件答案块样式基准，与 S5-W4 并轮）。判断场随页。\par}
\vspace{1.2mm}
\begin{multicols}{2}
\emergencystretch=1em

\tihao{1}\zu{单选} \tieside{简单}已知\(\overrightarrow{a}=(1,2,2)\)，则\(|\overrightarrow{a}|=\)\nobreak(\kongwei)

\lxoptq{A．\(1\)；}{B．\(2\)；}{C．\(3\)；}{D．\(4\)}

\begin{ansblock}[A试-3-jc-01]
% ans:A试-3-jc-01 C
\ansitem{1}{\(C\)（\(\sqrt{1+4+4}=3\)）．}
\end{ansblock}

\zhenhead{判断正误(正确的打\gou{},错误的打\cha{})}

\zhentib{(1)}{零向量与任意向量平行．}[\gou]

\zhentib{(2)}{方向相同或相反的向量是共线向量．}[\cha]

\begin{ansblock}[A试-3-jc-02]
\ansitem{(2)}{(1)\gou{}；(2)\cha{}（方向相同或相反即共线，但零向量与任意向量平行，命题漏零向量）．}
\end{ansblock}

\tihao{2}\zu{填空} \tieside{简单}向量\(\overrightarrow{a}\)，\(\overrightarrow{b}\)满足\(|\overrightarrow{a}|=3\)，\(|\overrightarrow{b}|=4\)，\(\overrightarrow{a}\perp\overrightarrow{b}\)，则\(|\overrightarrow{a}+\overrightarrow{b}|=\)\kongbai{}．

\begin{ansblock}[A试-3-jc-03]
% ans:A试-3-jc-03 5
\kbfill{5}
\ansitem{2}{\(5\)．}
\end{ansblock}

\end{multicols}

\end{document}
